\section{Train and Select a Model on Shared Compute}
\label{chap:training}

Training should begin only after the dataset interface has been frozen and a
small end-to-end smoke test has succeeded. The goal of the server workflow is
not merely to obtain a checkpoint; it is to preserve enough information to
reproduce how that checkpoint was produced.

\subsection{Freeze one training configuration}

Record the following before launching a long run:

\begin{itemize}
  \item dataset path and episode count;
  \item observation keys, shapes, and preprocessing;
  \item action representation and horizons;
  \item model architecture and parameter count;
  \item optimizer, learning rate, scheduler, batch size, and seed;
  \item train/validation split by episode;
  \item number of epochs and checkpoint interval;
  \item software revision and training device.
\end{itemize}

For a new model, first train for only a few batches and load the resulting
checkpoint in the policy runner. This verifies configuration composition,
tensor shapes, normalization, and serialization before consuming hours of GPU
time.

\subsection{Keep execution contexts explicit}

Distinguish the local experiment PC, server host, container, and Python
environment. Many path and import errors come from running a correct command in
the wrong context; for example, \filepath{/workspace} exists inside the project
container, not on the local robot computer.

\subsection{A general server workflow}

\begin{enumerate}
  \item Validate the dataset and record the source revision.
  \item Transfer only the required source, configuration, and data.
  \item Use a persistent session and a documented container/environment.
  \item Print resolved paths and run a short smoke test.
  \item Monitor the run and copy model candidates with their configurations
  back to the deployment computer.
\end{enumerate}

Use a persistent container or a documented image. Installing arbitrary packages
interactively until an import works makes the next run difficult to reproduce.

\subsection{Example: LIRMM training setup}

\begin{projectbox}[Internship example]
Training ran on \filepath{mic-gpu1} inside a persistent Docker container using
one NVIDIA A100 GPU. Data and configurations were transferred with
\filepath{rsync}; \filepath{tmux} protected the run from SSH disconnection. An
absolute dataset path was supplied explicitly, and the trained model was loaded
on the local CPU before robot testing. Exact server and training commands are in
\filepath{README_SERVER.md}.
\end{projectbox}

\subsection{What to monitor}

Monitor several signals rather than one number:

\begin{description}
  \item[Training loss:] confirms that the optimizer can fit the demonstrations.
  \item[Validation loss:] detects generalization degradation on held-out episodes.
  \item[Action error or samples:] helps reveal one problematic action component.
  \item[GPU use and step time:] reveals data-loader, memory, or device problems.
  \item[Closed-loop rollouts:] measure the behavior that ultimately matters.
\end{description}

The first Zarr load may be slow because chunks are decompressed and cached. A
pause after model construction is not proof that training has frozen; inspect
CPU, disk, GPU, and process activity before interrupting it.

\subsection{Select model candidates}

Save model candidates often enough to observe overfitting. Validation loss may
reach its minimum early, while closed-loop performance does not always rank
candidates in exactly the same order.

A sensible selection procedure is:

\begin{enumerate}
  \item reject runs with unstable or inconsistent training;
  \item choose a small set around the best validation region;
  \item benchmark inference time and load each checkpoint locally;
  \item test candidates in simulation or a no-motion real stage;
  \item choose one checkpoint before the reported evaluation trials.
\end{enumerate}

Do not change the selected model during the final evaluation. Development
trials and evaluation trials have different roles.

\subsection{Typical failures}

\begin{longtable}{@{}p{0.29\textwidth}p{0.65\textwidth}@{}}
  \toprule
  \textbf{Symptom} & \textbf{First checks} \\
  \midrule
  \endhead
  Dataset not found & Print the path inside the container; verify mount, transfer location, and Hydra-resolved value. \\
  CUDA out of memory & Check other processes, selected GPU, batch size, workers, image resolution, and model size. \\
  Training appears frozen & Inspect first-load caching, disk activity, workers, shared memory, and GPU utilization. \\
  Loss becomes NaN & Check normalization, invalid actions, learning rate, mixed precision, and exploding gradients. \\
  Validation rises early & Check split size, dataset diversity, augmentation, checkpoint frequency, and overfitting. \\
  Checkpoint fails locally & Check source revision, policy class, preprocessing, dependency compatibility, and EMA loading. \\
  \bottomrule
\end{longtable}

\begin{checkpointbox}[Training output]
For every candidate checkpoint, keep the resolved configuration, dataset name,
Git revision, epoch, validation metrics, parameter count, and a measured
inference latency on the deployment computer.
\end{checkpointbox}
